\cotefile{156-6}
\scanpage{2}\mspage{1}
\margin{\tikz[baseline=(g.base)]{\node[inner sep=2pt] (g) {GF \uline{VI}}; \draw (g.north west)--(g.south west)--(g.south east);}}
\begin{center}Analysis situs (deuxième mouture)\end{center}

% slanted
\margin{18 juin 86}Avant de décrire ce qu'est une \uline{contrée}, je vais décrire ce qui sera une ``contrée avec notion de multistrates'' -- la famille des multistrates \inserted{choisies} jouant un peu le rôle \unsure{de'|d'|des} une famille d'ouverts \deleted{\unsure{basiques|basique} pour} \inserted{engendrant} une topologie, ou une famille génératrice d'éléments d'un topos. Je vais donc commencer par

\textcircled{I}\quad ``Algèbre de figures'' ou ``\uline{\deleted{\unsure{a}}Atelier}''.

1.\quad Une \uline{algèbre de figures} implique avant tout trois types d'objets, les \uline{lieux}, les \uline{multistrates}, les \uline{figures}, \unsure{formant} trois ensembles
\[
\mathcal{L},\mathcal{M},\mathcal{F},\tag*{(1.1\unsure{.})}
\]
liés entre eux par diverses applications, et chacun muni de diverses structures. Ainsi, on a des applications canoniques injectives
\diagram{\begin{tikzcd}[ampersand replacement=\&,column sep=large]
{\text{(1.2)}} \& \mathcal{L} \arrow[r,hook,"\unsure{b)}"] \& \mathcal{M} \arrow[r,hook,"\unsure{a)}"] \& \mathcal{F}{,}
\end{tikzcd}}
que nous utiliserons souvent pour identifier un lieu à une multistrate particulière, et une multistrate à une figure particulière -- $\mathcal{L}$ à un \unsure{sous-ens.} de $\mathcal{M}$, $\mathcal{M}$ à un \unsure{sous-ens} de $\mathcal{F}$.

Il y a d'autre part deu\deleted{\unsure{x}}x\unsure{.} autres paires d'applications, que voici.
\scanpage{3}\mspage{2}
% slanted
\margin{b) se déduit de a) et de 1.2. a)}
\diagram{\begin{tikzcd}[ampersand replacement=\&,column sep=small,row sep=small,
execute at end picture={\draw (aa.north west) .. controls +(0,-.25) and +(.04,.16) .. (aa.south west) .. controls +(-.04,-.10) and +(-.16,.03) .. (bb.north west) .. controls +(.08,-.10) and +(-.02,.25) .. (bb.south west);}]
{\text{(1.3)}} \& |[alias=aa]| {\text{a)}} \& {\mathcal{F}\hookrightarrow\unsure{\mathcal{P}|\mathscr{P}}(\mathcal{M})} \& {F\mapsto\widetilde F=\operatorname{\text{Dépl}}(F)\subset\mathcal{M}} \\
\& |[alias=bb]| {\text{b)}} \& {\mathcal{F}\hookrightarrow\mathcal{P}(\mathcal{P}(\mathcal{M}))} \& {\text{plus précisément }\mathcal{F}\to\operatorname{Fig}(\mathcal{M})} \\
\& {\substack{\text{ne sert}\\\text{à rien}}}\arrow[u,bend left=24]\arrow[rr,bend right=13] \& \& {F\mapsto\deleted{\unsure{\sim}}\operatorname{Fig}_{\mathcal{M}}(F)\subset\mathcal{P}(\mathcal{M})}
\end{tikzcd}}
où $\operatorname{Fig}(\mathcal{M})$ désigne la partie de $\mathcal{P}(\mathcal{P}(\mathcal{M}))$ formée des figures ensemblistes dans $\mathcal{M}$. \deleted{\unsure{On}}

On peut considérer que la première application correspond à une relation entre $\mathcal{M}$ et $\mathcal{F}$, appelée relation d'incidence. Pour une figure $F$, $\widetilde F$ s'appelle l'ensemble des \uline{multistrates} \uline{incidentes}, ou le \uline{déployement} de la figure $F$. Si $X\in\mathcal{M}$, $F\in\mathcal{F}$, on dira que la multistrate $X$ est \uline{incidente} à la figure $F$, ou encore que c'est une \uline{strate de} \deleted{$F$}\inserted{\uline{la figure} $F$}, si $X\in\widetilde F$. D'autre part, tout\deleted{\unsure{e}} élément \inserted{$X$} de $\mathcal{M}$ (i.e. \deleted{\unsure{toute}} multistrate), regardé comme une figure par (1.2), admet un déployement $\widetilde X$, et \deleted{la multistrate} \deleted{\unsure{associée}} \deleted{\unsure{on}} on pose
\margin{$\unsure{\bullet|\cdot}$}
\[
\begin{aligned}
\operatorname{Fig}_{\mathcal{M}}(F)&=\{\widetilde X\mid X\in\widetilde F=\operatorname{Depl}(F)\}\subset\mathcal{P}(\mathcal{M})\\
\text{ie.}\quad \operatorname{Fig}_{\mathcal{M}}(F)&\in\mathcal{P}(\mathcal{P}(\mathcal{M}))
\end{aligned}\tag*{(1.4)}
\]
et il résultera des axiomes que c'est une figure ensembliste dans $\mathcal{M}$, \unsure{indexée} \unsure{fidélement|fidèlement} par l'\unsure{ens} $\widetilde F$ des strates de $F$.

\deleted{\unsure{En fait}} \deleted{\illegible{1}} En fait, $\mathcal{M}$ sera muni d'une relation d'ordre $\leq$, cf plus bas, et $\widetilde F\inserted{\subset\mathcal{M}}$ sera une partie fermée de \deleted{\unsure{la}} $\mathcal{M}$, et pour tout $X\in\widetilde F$, \deleted{\unsure{$X$}} on aura
\[
\widetilde X=\{Y\in\mathcal{M}\mid Y\leq\unsure{x|X}\}\tag*{(1.5\unsure{.})}
\]
\scanpage{4}\mspage{3}
À cause de cette interprétation, le passage de $\widetilde F\subset\mathcal{M}$ à $\operatorname{Fig}_{\mathcal{M}}(F)$ est à tel point évident, que \inserted{la considération de} cette figure ensembliste dans $\mathcal{M}$ me semble \unsure{rarement|vraiment} importante -- mais à voir.....

\deleted{Quand on \unsure{identi}}

On pourra utiliser les deux applications (1.3) pour identifier une figure, suivant les cas, \unsure{à} une partie de $\mathcal{M}$ [via $F\unsure{\mapstochar\relbar|\mapsto}\widetilde F$] ou à une figure ensembliste dans $\mathcal{M}$ [via $F\mapsto\operatorname{\text{Dépl}}(F)$]. Dans le premier cas, \unsure{comparant|composant} avec $\mathcal{M}\unsure{\lhook\joinrel\relbar|\hookrightarrow}\mathcal{F}$ (1.2)\deleted{\unsure{(}}, par la formule \unsure{1,5|1.5}. Cette dernière inclusion s'identifie à
\[
X\mathrel{\mapstochar\relbar\mapstochar}\widetilde X=\{Y\in\mathcal{M}\mid Y\leq X\}
\qquad \mathcal{M}\to\unsure{\mathcal{P}|\mathscr{P}}_{\mathrm{ferm}}(\mathcal{M}).
\]

Donc il faudra être prudent \deleted{\unsure{alors}} \inserted{\unsure{avec}} l'identification dont il a été question, des\deleted{\unsure{une|un}} $X\in\mathcal{M}$ \unsure{ou} multistrates à des figures particulières, quand celles-ci sont \unsure{interprétées} comme des parties de $\mathcal{M}$: La figure associée à $X\in\mathcal{M}$ n'est nullement $\{X\}$, mais $\widetilde X\unsure{\rbrace|,|;}$ \deleted{\unsure{$F$|l'}} \unsure{ens} des multistrates majorées par $\unsure{x|X}$ pour $\leq$ (ou encore, des sous-multistrates, ou multistrates incidentes à $X$, ou encore des \deleted{\unsure{$F$}} \uline{strates} de $x$).
\scanpage{5}\mspage{4}
D'autre part, il y a deux \inserted{\uline{autres}} applications
\[
\left\{\begin{aligned}
\text{a)}\quad &\unsure{\mathcal{F}|\mathcal{E}}\longrightarrow\unsure{\mathcal{P}|\mathscr{P}}(\deleted{\unsure{\mathcal{F}}}\mathcal{M})&
 F&\longmapsto\operatorname{Omb}_{\unsure{\mathcal{M}|m}}(F)\\
\text{b)}\quad &\mathcal{F}\hookrightarrow\mathcal{P}\mathcal{P}(\deleted{\unsure{\mathcal{F}}}\mathcal{M})&
 \text{et }\unsure{=}\quad\mathcal{F}&\longrightarrow\operatorname{Fig}(\mathcal{M}\unsure{]|)}\subset\mathcal{P}\mathcal{P}\unsure{\mathcal{M}|(\mathcal{M})}\\
&& F&\mapsto\deleted{\unsure{O}}\operatorname{Multomb}_{\mathcal{M}}(F),
\end{aligned}\right.\tag*{(1.\deleted{\unsure{4}}6)}
\]
où $\operatorname{Omb}_{\mathcal{M}}(F)$ s'appelle l'\uline{\unsure{ogombre|ombre}} \uline{de la} figure $F$ sur $\mathcal{M}$, ou la \uline{grande} \uline{ombre} \uline{de} $\mathcal{F}$. \deleted{\unsure{La}}Une multistrate \inserted{$X$} qui est dans la grande ombre d'une figure $F$ est dite ``\uline{raffiner}'' \uline{la figure} $F$ (c'est cette fois déduit d'une relation d'ordre $\ll$ sur $\mathbb{F}$, le raffinement, par la condition $X\ll F$). Il n'est pas clair que cette application $\mathcal{F}\unsure{\relbar|\to}\mathcal{P}(\mathcal{M})$ soit injective, ie. que $\mathcal{F}$ s'identifie à un \unsure{ens} de parties de $\mathcal{M}$, via cette application. Par contre il en sera ainsi de (1.6) b), qui est déduit de a) via \deleted{\unsure{$F$}}la \unsure{formation|fonction} de $\widetilde F$, comme (1.3 b) est déduit de (1.3. a):
\[
\operatorname{Multomb}(F)=\{X\in\widetilde F\mid\operatorname{Omb}X\},\tag*{(1.\deleted{\unsure{5|6}}7)\quad\deleted{\unsure{(4.5)|(1.5)}}}
\]
alors que $\operatorname{Omb}(X)$ s'explicite par
\[
\unsure{\operatorname{Multomb}|\operatorname{MultOmb}}(X)=\{Y\in\deleted{\unsure{\mathcal{F}}}\mathcal{M}\mid Y\ll X\},\tag*{(1.\deleted{\unsure{6}}8)}
\]
analogue à (1.5), où cette fois on a utilisé la relation \inserted{d'ordre} $\ll$ ald $\leq$ -- \deleted{\unsure{C'est là}}
\scanpage{6}\mspage{5}
\deleted{une figure a}

Pour \deleted{toute} figure $F$, $\operatorname{Multomb}(F)$ \inserted{est appelé} \deleted{\unsure{la|de}} la \uline{grande} \uline{multiombre} de $F$, ou la $\mathcal{M}$-multiombre de $\mathbb{F}$. \deleted{\unsure{On}} C'est en fait une figure ensembliste dans $\mathcal{M}$\inserted{indexée \unsure{fidèlement|fidélement} par $\widetilde F$ (le déployement de $F$)}, et dont la connaissance implique celle de $F$.

\inserted{Jusqu'à présent il n'était pas question de $\mathcal{L}$.} Utilisant \deleted{d'autre part} \inserted{\unsure{maintent|maintnt|maintenant}} l'application
\[
\mathcal{L}\hookrightarrow\mathcal{M},
\]
et les applications images inverses correspondantes
\[
\begin{aligned}
\unsure{\mathcal{P}|\mathscr{P}}(\mathcal{M})&\longrightarrow\mathcal{P}(\mathcal{L})\\
\operatorname{Fig}(\mathcal{M})&\unsure{\rightharpoonup|\longrightarrow|\relbar}\operatorname{Fig}(\mathcal{L})
\end{aligned}
\]
on \unsure{trouve|trouvera}, en composant avec 1.6 a)b)
\[
\left\{\begin{aligned}
\text{a)}\quad \mathcal{F}&\mathrel{\deleted{\unsure{\lhook|\leftarrow}}\!\longrightarrow}\mathcal{P}(\mathcal{L})&
 F&\longmapsto\operatorname{omb}(F)\\
\text{b)}\quad \mathcal{F}&\mathrel{\deleted{\lhook}\!\longrightarrow}\operatorname{Fig}(\mathcal{L})&
 F&\longmapsto\operatorname{multomb}(F),
\end{aligned}\right.\tag*{(1.9)}
\]
dont la première n'a aucune raison d'être injective, la seconde par contre a une tendance à l'être\unsure{$\times$|*|.} \deleted{\unsure{on ne|a priori} rien} \inserted{On} appelle $\operatorname{omb}(F)$ la \uline{petite ombre} de $F$, ou simplement son ombre quand une confusion semble exclue, et \unsure{de} $\unsure{\overline m|=}$ pour $\operatorname{multomb}$, la (\uline{petite}) \uline{multiombre de la} figure $F$.
\scanpage{7}\mspage{6}
Ainsi, l'ombre d'une figure est formée des lieux qui la raffinent -- on dit aussi que ce sont les \uline{lieux ``de la figure}'' (à ne pas confondre, avec les strates de la figure qui sont des lieux). La \deleted{\unsure{om}}multiombre est donc \unsure{formé|formée} des ombres de ses strates, elle est indexée par le déployement $\widetilde F$, ensemble des strates de $F$, la relation d'\deleted{\unsure{d'ordre|d'inclu}}\unsure{inclusion} de la figure ensembliste n'étant autre que la relation d'\unsure{incidences|incidence} entre strates de $F$, ie. la relation \unsure{de|dans} $\widetilde F$ induite par $\leq$ \unsure{de|dans} $\mathcal{M}$ (ou par $\ll$ au choix, \unsure{ça|sa} revient au $\unsure{\overline m|=}$). Alors qu'une figure ne se réconstitue pratiquement jamais à partir de son ombre, elle se réconstitue \unsure{pourtant} \deleted{\unsure{à}}souvent (cas des ``\deleted{\unsure{a}}Ateliers ensemblistes'') \deleted{\unsure{de}}à partir de sa multiombre. Nous y reviendrons.

\unsure{\S|S} Il y a lieu de parler des relations d'ordre canoniques
\scanpage{8}\mspage{7}
\[
\begin{aligned}
\leq,\quad\ll,\quad\unsure{\preccurlyeq|\preceq}&\qquad\text{relations d'ordre}\\
&\qquad\text{de }\mathcal{F}
\end{aligned}\tag*{(1.10).}
\]
et les relations induites dans $\mathcal{M}$
\[
\leq,\quad\ll\quad\text{dans }\mathcal{M}\tag*{(1.11)}
\]
% slanted
\margin{\uline{NB} les relations d'ordre induites par $\leq$, $\ll$ sur $\mathcal{L}$ sont discrètes}
[\unsure{nb|NB} la relation d'ordre $\preccurlyeq$ induite dans $\mathcal{M}$ est moins importante, comme structure spécifique dans $\mathcal{M}$].

% slanted
\margin{\uline{NB} $\unsure{G|C}\leq F\Longrightarrow G\ll F$ une sous-figure de $F$ raffine $F$}
\diagram{\(\begin{array}{@{}l@{\qquad}p{.64\linewidth}@{}}
\deleted{F}G\deleted{\unsure{<|\ll}}\leq\deleted{G}F & $\deleted{F}G$ est dit une \uline{sous}-\uline{figure} de $\deleted{G}F$\\
& (Si $\deleted{F}G\inserted{=X}\in\mathcal{M}$, c'est donc une \uline{strate} de $\deleted{G}F$)\\[.45em]
G\ll F & $G$ est dit un \uline{raffinement} de $F$\\
& (si $G=\unsure{x|X}\in\mathcal{M}$, c'est donc un \unsure{élément|élt} de la grande o\deleted{\unsure{o|m}}mbre de $F$)\\[.45em]
G\preccurlyeq F & $G$ est dit une \uline{subdivision} de $F$,\\
& on dit que la figure $G$ subdivise la figure $F$.
\end{array}\)}

Je note enfin deu\deleted{\unsure{x}}x autres relations importantes \deleted{\unsure{dans}} entre figures:
\[
\left\{\begin{aligned}
F\text{ et }G&\quad\text{\uline{compatibles}}\\
F\text{ et }G&\quad\text{\uline{disjoints}}\qquad,
\end{aligned}\right.\tag*{(1.12)}
\]
qui l'une et l'autre peuvent s'exprimer en termes de la relation d'ordre $\leq$:

$F$ et $G$ compatibles: ce sont des sous figures d'une $\unsure{\overline m|=}$ figure, \unsure{i..e|i.e} $\{F,G\}$ majoré pour $\leq$;

$F$ et $G$ \unsure{disjoint|disjoints}: $F$ et $G$ compatibles, et leur
\scanpage{9}\mspage{8}
$\inf$ dans $\mathcal{F}$ est $\varnothing_{\mathcal{F}}$, le plus petit élément de $\mathcal{F}$ ou \uline{figure vide}. Ainsi
\[
F\text{ et }\unsure{G|C}\text{ disjoints}\quad\Longrightarrow\quad F\text{ et }G\text{ compatibles}.
\]

On a, par les axiomes:
\[
\begin{aligned}
&&&\deleted{\text{tout}}\quad\inserted{\operatorname{omb}(F)\text{ et }\operatorname{omb}(G)}\\
 F\text{ et }G\text{ disjoints}&\quad\Longleftrightarrow\quad&&\deleted{\operatorname{omb}(F)\cap\operatorname{omb}(G)}\\
&&&\text{sont ``disjoints'' au sens}\\
&&&\text{suivant: ils sont \unsure{disjoints|disjoint}}\\
&&&\text{ensemblistement \unsure{i.e.|ie.}}\\
&&&\operatorname{omb}(F)\cap\unsure{\operatorname{om}|\operatorname{omb}}(G)=\varnothing,\ \text{et}\\
&&&\text{tout \unsure{élément|élém\textsuperscript{t}} de }\operatorname{omb}(F)\text{ est}\\
&&&\text{compatible avec tout \unsure{élt|élément}}\\
&&&\text{de }\operatorname{om}(G)
\end{aligned}\tag*{(1.13)}
\]
-- i.e. deux figures sont disjointes ssi elles n'ont pas de \unsure{lieu|lieux} communs, et \unsure{tout|$\forall$} lieu de l'une est \deleted{disjoint} \deleted{\unsure{com}} \inserted{disjoint} (au sens de \deleted{la} l'\deleted{théorie}atelier) de tout lieu de l'autre;

Cela montre que les relations de \deleted{\unsure{disj}}compatibilité et de disjonction sont importantes aussi dans les \deleted{\unsure{en de|en des}} \inserted{sous-}ensembles $\mathcal{M},\mathcal{L}$ de $\mathcal{F}$. Dans le cas des ateli\deleted{\unsure{er|ers}}ers ensemblistes, ces relations sont ``triviales'' sur $\mathcal{L}$: deux lieux sont tjrs \deleted{disjoints} \inserted{compatibles}, et disjoints ssi ils sont distincts. Par contre, même dans le cas ensembliste et $\unsure{\overline m|=}$ modéré, la relation
